
% Lecture Template for ME3001-001-Tristan Hill - Spring 2017 - Fall 2017 - Fall 2020
% Mechanical Engineering Analysis with MATLAB
% Module 1 - Introduction 
% Topic 5 - MATLAB Review - Loops

% Document settings

%\documentclass{beamer}                  % for presentation ?
\documentclass[handout]{beamer}  % for handout ?
\usepackage{/home/thill/Documents/lectures/analysis_lectures/analysis_lectures}

\newcommand{\MNUM}{1\hspace{2mm}} % Module number
\newcommand{\TNUM}{5\hspace{2mm}} % Topic number 
\newcommand{\moduletitle}{Introduction and MATLAB Review} % Titles and Stuff
\newcommand{\topictitle}{MATLAB Loops} 

\newcommand{\sectiontitleI}{What is Looping?} % More Titles and Stuff
\newcommand{\sectiontitleII}{Types of Loops}
\newcommand{\sectiontitleIII}{MATLAB Examples}


\newcommand{\btVFill}{\vskip0pt plus 1filll}

% custom box
\newsavebox{\mybox}

\author{ME3001 - Mechanical Engineering Analysis} 
\title{Lecture Module - \moduletitle}
\date{Mechanical Engineering\vspc Tennessee Technological University}

\begin{document}
	
	\lstset{language=MATLAB,basicstyle=\ttfamily\small,showstringspaces=false}
	
	\frame{\titlepage \center\begin{framed}\Large \textbf{Topic \TNUM - \topictitle}\end{framed} \vspace{5mm}}
	
	% Section 0: Outline
	\frame{
		\large \textbf{Topic \TNUM - \topictitle} \vspace{3mm}\\
		
		\begin{itemize}
			
			\item \sectiontitleI    \vspc % Section I
			\item \sectiontitleII 	\vspc % Section II
			\item \sectiontitleIII 	\vspc %Section III
			%\item \sectiontitleIV 	\vspc %Section IV
			
		\end{itemize}
		
	}



\section{\sectiontitleI}

\frame{
  \frametitle{\sectiontitleI}
  
  
Repeated Code


You may have you noticed in your programs that there was a lot of \PR Repeated Code \BK? Is this OK?\vspc

	\begin{itemize}
		\item 
		\item \vspace{3mm}
	\end{itemize}
	\vspace{3mm}

What can we do about it?  
            			
Introduce a new \color{blue}Control Structure \color{black} ... 
		
}


\frame{
\frametitle{A New Control Structure}
	

A loop is a fundamental {\bf control structure} in sequential programming in which one or more commands is executed in repetition controlled by a logical statement or counter. - TH  
\vspace{5mm}\\
... In computer programming, a loop is a sequence of instructions that is repeated until a certain condition is reached.  - Web

	\begin{itemize}
		\item 
		\item \vspace{3mm}
	\end{itemize}
	\vspace{3mm}

}


\frame{
\frametitle{Why use Loops? }


         \Large Why are loops useful in computer programming? \vspc
	\large
	\begin{itemize}
		\item Reduce problems associated with {it repeated code} \vspc
		\item \vspace{3mm}
		\item \vspace{3mm}
	\end{itemize}

}


\section{\sectiontitleII}
\frame{
\frametitle{Many Different Types}

\large There are many different types of loops in computer programming. This is not an inclusive list.\vspc
\begin{itemize}

	\item \BL WHILE \BK loops
	\item \BL FOR \BK loops
	\item {\it nested} loops
	\item DO... WHILE loops
\end{itemize}

\large Loops can be characterized as {\it entry-controlled} or {\it exit-controlled} loops.

}

\frame{
\frametitle{The WHILE Loop}

 The WHILE Loop is  the simplest loop in computer programming.\vspc
\begin{itemize}

	\item {\bf block of code} is repeated as long as \textbf{\PR condition} \BK is\textbf{ \GR true } \BK \vspc
	\item {\it entry-controlled} loop structure \vspc
	\item \vspace{3mm}
	\item  \vspace{3mm}

\end{itemize}

}

\frame[containsverbatim]{
\frametitle{Structure of a WHILE Loop}

%It is useful to visualize the structure in a {\bf program flowchart} \vspcc
%%\begin{multicols}{2}
%\includegraphics[scale=0.25]{lecture1_fig1.png}\\	
%
\begin{framed}
\begin{lstlisting}
i=1; %Initialize the Counter

while(i<10)

	% Do Stuff Here

	i=i+1; % Increment the Counter

end	
\end{lstlisting}
\end{framed}
%
%%\end{multicols}
}


%Section 3: MATLAB Examples
\section{\sectiontitleIII}
% references is not a section for now, for looks and it would be a waste of space
\frame{

\frametitle{References}

\begin{itemize}
	\item Your MATLAB textbook - Chapter 5 - Looping Statements
\end{itemize}

}
\end{document}









 